\documentclass[journal,a4paper]{IEEEtran}

\usepackage{mir-template}

\title{Preparation of Papers for Machine Intelligence Research}
\author{First Author\textsuperscript{1}\qquad Second Author\textsuperscript{1,2}}
\miraffiliation{%
\textsuperscript{1}Editorial Office of Machine Intelligence Research, Institute of Automation, CAS, Beijing 100190, China\\[3pt]
\textsuperscript{2}Department of Academic Journals, Institute of Automation, CAS, Beijing 100190, China}

\begin{document}

\begin{abstract}
An abstract should be a concise summary with no more than 500 words of the significant items in the paper, including the background, methods, results, and conclusions. Define all nonstandard symbols, abbreviations and acronyms used in the abstract. Do not cite references in the abstract.
\end{abstract}

\begin{IEEEkeywords}
Keyword 1, keyword 2, keyword 3, keyword 4, keyword 5.
\end{IEEEkeywords}

\maketitle

\section{Introduction}

\textit{Machine Intelligence Research} (MIR) is published by Springer, and sponsored by Institute of Automation, Chinese Academy of Sciences. MIR publishes papers on original theoretical and experimental research and emerging technologies in intelligent science, and strives to bridge the gap between theoretical research and practical applications. Topics include artificial intelligence basic theory, machine learning, pattern recognition \& computer vision, natural language processing, knowledge management \& data mining, brian-inspired intelligence, intelligent robotics, distributed artificial intelligence \& multi-agent systems, and innovative AI applications.

The instruction aims to help you preparing papers for MIR. Download the template files from the website at www.mi-research.net (to be live on 1/1/2021) and revise your manuscript according to the requirements. Use Word software for formatting.

\section{Procedure}

\subsection{Submission of papers}

Please log in MIR website at www.mi-research.net to submit your papers. Papers must be written in English. All submissions are acknowledged. DOC, PDF or PS files that are ready for printing can be accepted.

Submitted articles may be of the four basic types:

\textbf{Review:} Extensive reviews of established or emerging research topics or application areas ( $>$ 10 printed pages).

\textbf{Research article:} Detailed discussion involving new research, applications or developments ($>$10 printed pages).

\textbf{Correspondence:} Letters to the Editor about the journal or to authors commenting on previously published papers. In the latter case, the Editor will give the authors an opportunity to respond (at least 1 printed page).

\textbf{Perspective:} Perspective (at least 1 printed page) highlights recent exciting research, but do not primarily discuss the author's own work. The authors may provide context for the findings within a field or explain potential interdisciplinary significance.

\subsection{Review stage}

All manuscripts submitted to MIR are subject to peer review. Members of the Editorial Board choose the referees and make the final decision of acceptance or rejection. Most manuscripts are returned to authors for revisions recommended by the editor and referees. Thus, it will typically take some months for a paper to be finally accepted.

\subsection{Final stage}

Along with the Notice of Acceptance, a Transfer of Copyright Form will be attached which should be signed by the author so as to transfer the copyright of the written work to the journal and permit the following publication process. The Transfer of Copyright Form should be properly completed and signed after a paper has been accepted.

After your paper has been accepted, send complete contact information of all authors, including institutes, full mailing addresses, telephone numbers, ORCID iD of authors, and e-mail addresses to MIR when submitting the final version. In addition, designate one author as the ``corresponding author''.

For typeset papers, proofs are emailed to all authors, and the proofed version should be returned before the due time. Corrections should be restricted to typesetting errors. Otherwise, the publication process will be delayed because of repetitive check and review. Any queries from editors should be answered in full.

\section{Manuscript preparation}

Complete manuscripts should include the following parts: title, authors' names, authors' affiliations, abstract, keywords, text, acknowledgments (if necessary), appendixes (if necessary), collected references listed corresponding to their citation order, short biography and the photograph of every author (above 600 dpi), tables and figures (above 600 dpi).

\subsection{References}

All publications cited in the text should be presented in the reference list at the end of the manuscript. References must be listed in the order they were cited (numerical order). Number citations consecutively in square brackets\cite{lecun2015deep}. The sentence punctuation follows the brackets\cite{vaswani2017attention}. Multiple references \cite{vaswani2017attention,he2016deep} are numbered with one brackets, such as \cite{lecun2015deep,vaswani2017attention,he2016deep}, \cite{vaswani2017attention,he2016deep}, \nocite{bishop2006pattern}\cite{he2016deep,radford2021learning,deng2009imagenet,mnih2015human}, etc. When citing a section in a book, please give the relevant page numbers\cite{bishop2006pattern}. It is not necessary to mention the author(s) of the reference unless it is relevant to your text. In sentences, refer simply to the reference number, as ``in \cite{radford2021learning}''. Do not use ``Ref. \cite{radford2021learning}'' or ``reference \cite{radford2021learning}'' except at the beginning of a sentence: ``Reference \cite{radford2021learning} shows \textellipsis''.

Please note that the references at the end of this document are in the preferred referencing style. The names of all authors should be given in the references unless authors are more than six persons. If there are more than six authors, you may use ``et al''. Use a space after authors' initials. Papers that have been submitted for publication but have not been accepted should not be cited. Preprints and online resources should include a stable URL or DOI whenever possible\cite{kingma2015adam}. Such a reference may be updated at the proof stage if the referenced paper has been published by then. References should not contain abbreviations or acronyms for the names of Journals, Conferences, Laboratories or Companies; spell their full names.

\subsection{Tables and figures}

As shown in Table 1, place table titles above the tables. Do not abbreviate ``Table''. Statements that serve as captions for the entire table do not need footnote letters.

Number figures in accord with the order they appeared in the text and make sure that every figure is cited.

Every figure must have a caption that is complete and intelligible in itself without reference to the text.

Because MIR will do the final formatting for your paper, you do not need to put figures and tables at the top and bottom of each column. Instead, all figures, figure captions, and tables can be put at the end of the paper. Large figures and tables may span across two columns. If one figure has two parts, mark them as (a) and (b), rather than (1) and (2). Please make sure that the figures and tables you mentioned in the text actually exist. Place the figure captions below the figure, and captions should be separated from the figures. Use the abbreviation ``Fig.'' even at the beginning of a sentence, e.g., Fig. 1, Figs. 2-4, Figs. 5 and 6.

The preferred format is encapsulated postscript (.jpg, .eps) for figures with a minimum resolution of 600 dpi (dots per inch) which can be created by Matlab directly, or transformed by Photoshop software. The description style of the figures is specified as ``Times New Roman 8pt-12pt''.

Do not use color unless it is necessary to interpret your figures or add scientific information which is unavailable in an equivalent monochrome version. The author is required to afford additional cost for color printing, but there is no need to pay extra fees for electronic color copies. Consult the MIR Editorial Office for details.

Use words rather than symbols in axis labels to avoid confusion. For example, write ``Magnetization,'' or ``Magnetization M,'' rather than ``M''. Put units in parentheses. Do not label axes only with units. As in Fig. 1, for example, write ``Magnetization (A/m)'', not ``A/m''. Do not label axes with a ratio of quantities and units. For example, write ``Temperature (K)'', not ``Temperature/K.''

Clarify multipliers. For example, write ``Magnetization (kA/m)'' or ``Magnetization (103 A/m)'' rather than ``Magnetization (A/m) $\times$ 1000'' because the reader may not know whether the top axis label in Fig. 1 means 16000 A/m or 0.016 A/m.

\begin{figure}[!bp]
  \centering
  \includegraphics[width=\samplefigonewidth]{fig1.png}
  \caption{Magnetization as a function of applied field. Note that ``Fig.'' is abbreviated. There is a period after the figure number, followed by two spaces. It is good practice to explain the significance of the figure in the caption.}
\end{figure}

\begin{figure}[!b]
  \centering
  \begin{minipage}{\samplefigtwowidtha}
    \centering
    \includegraphics[width=\linewidth]{fig2a.png}
    {\mircaptionfont (a)\quad Scale 1\par}
  \end{minipage}

  \vspace{3pt}
  \begin{minipage}{\samplefigtwowidthb}
    \centering
    \includegraphics[width=\linewidth]{fig2b.png}
    {\mircaptionfont (b)\quad Scale 2\par}
  \end{minipage}
  \caption{Real value (solid) and estimated value (dotted)}
\end{figure}

\begin{table*}[!t]
  \centering
  \caption{Simulation result comparison between CSU\_Yunlu and other RoboCup teams}
  \begin{tabular*}{\mirwidetablewidth}{@{\extracolsep{\fill}}c c c c >{\centering\arraybackslash}p{0.24\textwidth}@{}}
    \toprule
    OPP & CSU:OPP & Group decision (\%) & Unit decision (\%) & Group decision / Unit decision \\
    \midrule
    Cyberoos2001 & 3:0 & 71.5 & 28.5 & 2.51\\
    FCPortugal2001 & 1:0 & 68.4 & 31.6 & 2.16\\
    Gemini & 26:0 & 59.7 & 40.3 & 1.48\\
    Harmony & 3:0 & 69.9 & 30.1 & 2.32\\
    Lazarus & 11:0 & 57.3 & 42.7 & 1.34\\
    MRB & 2:0 & 63.2 & 32.8 & 1.93\\
    SBCe & 4:1 & 65.8 & 34.2 & 1.92\\
    UvA\_Trilearn\_2001 & 1:0 & 54.9 & 44.1 & 1.24\\
    UTUtd & 10:0 & 70.7 & 29.3 & 2.41\\
    WrightEagle2001 & 3:1 & 66.2 & 33.8 & 1.96\\
    Average & & 64.8 & 35.2 & 1.84\\
    \bottomrule
\end{tabular*}

\end{table*}

\begin{table}[!t]
  \centering
  \caption{Experimental results of close-up retrieval}
  \begin{tabular*}{\mircolumntablewidth}{@{\extracolsep{\fill}}ccc@{}}
    \toprule
    Video clips & \textit{Recall (\%)} & \textit{Precision (\%)}\\
    \midrule
    Movie \#1 & 90.63 & 81.69\\
    Movie \#2 & 92.86 & 78.00\\
    TREC2001 NAD32 & 86.67 & 89.66\\
    TREC2001 NAD55 & 95.00 & 87.69\\
    Sports \#1 & 93.42 & 92.21\\
    News \#1 & 94.74 & 94.74\\
    Average & 92.22 & 87.33\\
    \bottomrule
\end{tabular*}

\end{table}

\subsection{Equations}

Number equations consecutively throughout the text with Arabic numbers in parentheses: (1), (2), (3), etc. In appendixes, use numerical sequence: (A1), (A2), (A3), etc.

Clearly define the symbols in equations. Italicize symbols (T might refer to temperature, as well as the unit of tesla). Refer to ``(1)'', not ``Eq. (1)'' or ``equation (1)'', except at the beginning of a sentence: ``Equation (1) is ...''.

\begingroup
\small
\begin{equation}
\left(
\begin{array}{cc}
A_{cl}^{T}P+PA_{cl} & PB_1\\
B_1^{T}P & -I
\end{array}
\right)<0
\end{equation}
\endgroup

\begingroup
\small
\begin{equation}
\begin{aligned}
g_i(t)=i\!\int_t^\infty e^{(A_1-SP_1)^T(r-t)}
&\left[P_1A_1x^{(i-1)}(r-\tau)\right.\\
&+A_1^TP_1x^{(i-1)}(r+\tau)\\
&\left.+A_1^Tg_{i-1}(r+\tau)\right]\,dr
\end{aligned}
\end{equation}
\endgroup

\subsection{Abbreviations and acronyms}

Define abbreviations and acronyms when they first appeared in the text even they have already been defined in the abstract or keywords. There is no need to define abbreviations such as IEEE, SI, etc.. Do not add space between alphabets in abbreviated words: write ``C.N.R.S.'', rather than ``C. N. R. S.''.

\section{Publication principles}

To avoid duplicate publication, manuscripts submitted to MIR, except for rejection ones, are not allowed to submit or publish elsewhere in English or any other language unless getting a written permission from MIR Editorial Office.

Do not submit a paper you have submitted or published elsewhere. Do not publish ``preliminary'' data or results. It is the obligation of the authors to cite prior relevant work. The author who submits the paper should obtain publication agreement from all coauthors and any permission required from sponsors. Each author should make direct contribution to the paper, and it is authors, not MIR, who are responsible for the content.

\section{Conclusions}

Typical functions of the conclusion include: 1) summing up, 2) a statement of conclusions, 3) a statement of recommendations, and 4) a graceful termination. Any one of these, or any combination, may be appropriate for a particular paper. Some papers do not need a separate conclusion part, particularly if the conclusions have already been stated in the introduction.

\section*{Appendix}

Appendixes, if needed, appear before the acknowledgment.

\section*{Acknowledgements}

In general, limit acknowledgments to those who helped directly in the research itself or during discussions on the subject of the research. Financial support of all kinds acknowledgments are placed here.

Use the singular heading even if you have many acknowledgments. Avoid expressions such as ``One of us (S.B.A.) would like to thank ...''. Instead, write ``First A. Author thanks ...'' .

Sponsor and financial support acknowledgments, if any, can be stated here. This work was supported by National Natural Science Foundation of China (No. xxxxxxxx) .

\nocite{krizhevsky2012imagenet,paszke2019pytorch,goodfellow2014generative}
\bibliographystyle{mir}
\bibliography{reference}

\begin{IEEEbiography}[{
\includegraphics[width=\mirauthorphotowidth]{author-first.png}}]{First Author}
and the other authors may include biographies at the end of the paper. Locate the photographs of authors on the left of the biographies. The author's educational background is listed. The degrees should be listed with type of degree in what field, which institution, city, state or country, and the year degree was earned. Listing military and work experience, including summer and fellowship jobs at the end of this paragraph. The current job must have a location.

The second paragraph uses the pronoun of the person (he or she) and not the author's last name. Information concerning previous publications may be included. Try not to list more than three books or published articles. The format for listing publishers of a book within the biography is: title of book (city, state: publisher name, year) similar to a reference. Current and previous research interests should be stated.

The third paragraph begins with the author's title and last name (e.g., Dr. Smith, Prof. Jones, Mr. Kajor, Ms. Hunter). List any memberships in professional societies. Finally, list any awards and work for committees and publications. If a photograph (above 600 dpi) is provided, the biography will be indented around it. The photograph is placed at the top left of the biography. Personal hobbies will be deleted from the biography. Please see the following example for more details.

E-mail: mir@ia.ac.cn (Corresponding author)

ORCID iD: 0000-0000-0000-0000
\end{IEEEbiography}

\begin{IEEEbiography}[{
\includegraphics[width=\mirauthorphotowidthsmall]{author-second.png}}]{Second Author}
received the B.Sc. and M.Sc. degrees in mechanical engineering from *** University, China in 1977 and 1984, respectively, and the Ph.D. degree in computing from *** University, UK in 1992. In 1994, he was a faculty member at *** University, China and in 1996 at *** University, USA. Currently, he is a professor in Department of Information System Engineering at *** University, China. Prof. Author received research award from Science Foundation, and the Best Paper Award of the IS International Conference in 2000 and 2006, respectively. He is a member of SICE, IEE and IEEE.

His research interests include robotics, feedback control systems, and control theory.

E-mail: \href{mailto:mir@ia.ac.cn}{mir@ia.ac.cn}

ORCID iD: 0000-0000-0000-0000
\end{IEEEbiography}

\end{document}
